\chapter{Introducción}

\section{Objetivo de la práctica}

El objetivo principal de esta práctica es utilizar \textit{OpenSSL} para realizar 
operaciones criptográficas utilizando:

\begin{itemize}[label=\textbullet]
    \item Algoritmos de resumen (\textit{hash})
    \item Algoritmos de cifrado simétrico
    \item Diferentes modos de operación
    \item Derivación de claves a partir de contraseñas
\end{itemize}

\section{Enunciado entero:}
\subsection{Generación y comprobación de Resúmenes}
Una vez entendido el empleo del comando openssl-dgst y sus opciones básicas (-help, -list):
\begin{enumerate}
    \item Crear un archivo de texto legible de entre 150 y 200 caracteres.
    \item Aplicar un mínimo de CINCO algoritmos de resumen distintos (md5, SHA-1, SHA-2 y
    Whirlpool son obligatorios) sobre ese archivo de texto, verificar que el tamaño coincide con
    la descripción del algoritmo de resumen empleado y comprobar cómo varían los resúmenes
    obtenidos con el mismo algoritmo ante mínimas modificaciones del fichero (cambiando un
    solo carácter, por ejemplo).
    \item Alternar entre salida hexadecimal, hexadecimal con : y binaria
\end{enumerate}
\subsection{Cifrado Simétrico de documentos}
Una vez entendido el empleo del comando openssl-enc para cifrar y descifrar, sus opciones básicas
(-help y -list) con diferentes algoritmos y modos de operación, el empleo de 
ficheros binarios y BASE64, la derivación de claves a partir de contraseñas 
y sal, detallada en el estándar PKCS \#5 (PBKDF1 y PBKDF2) y su aplicación a 
las claves de cifrado simétrico y vectores de inicialización:
\begin{enumerate}
    \item Crear un archivo de texto legible de pequeño tamaño – entre 31 y 81 caracteres (número impar).
    \item Cifrarlo (con salida en modo binario) con SIETE algoritmos simétricos (AES y TDES
    obligatorios, en modo bloque y flujo y dos cifradores de flujo -rc4 y –chacha20 o similar-).
    \item Descifrarlos y comprobar el resultado
    \item Explicar el tamaño de los diferentes ficheros cifrados en virtud del tamaño de bloque del cifrador
    (o no, si se cifra en flujo), y sabiendo que el empleo de sal añade 16 bits de más al inicio del
    fichero cifrado –Salted\_\_XXXXXXXX-.
    \item Cifrar el archivo con aes-256-cbc y contraseña y 
    descifrarlo NO con dicha contraseña, sino con su conjunto 
    equivalente de clave (key) y vector de inicialización (iv). 
    Para ello habrá que usar “-p” en el cifrado y eliminar los 16 
    bits iniciales del fichero cifrado que contienen la sal aleatoria 
    utilizada por el comando de cifrado (con el comando dd bs=1 skip=16 if= of= o 
    similar) antes de descifrarlo con -k y -iv, poniendo sus valores 
    mostrados al cifrarlo por la opción “-p”.
\end{enumerate}

\subsection{Aplicación: Demostración de la peligrosidad del modo de operación ECB}
En este apartado demostraremos que el modo de operación “ecb” con cifradores simétricos es muy
peligroso, repitiendo el ejemplo de Wikipedia con el pingüino Tux, utilizando una imagen pequeña
con colores sólidos como entrada. Para ello:
\begin{enumerate}
    \item Escogeremos una imagen de entre 20 y 50 Kb de colores sólidos (gif, png, etc.)
    \item Con la utilidad “convert” del paquete ImageMagick o similar la convertiremos a formato PGM
    (ver \url{https://en.wikipedia.org/wiki/Netpbm}).
    \item Separaremos la cabecera (3 primeras líneas de texto) y el cuerpo (binario) de la imagen
    \item Cifraremos el cuerpo con el cifrador de bloques más potente en la actualidad (AES-256) en modo
    ECB y una clave o contraseña aleatorias, si se usa contraseña, realizar el cifrado sin sal (-nosalt).
    \item Crearemos la nueva imagen con la cabecera anterior y el cuerpo cifrado
    \item Convertiremos la nueva imagen al formato original con Imagemagick “convert”
    \item Verificaremos que la nueva imagen puede ser “vista” de la misma forma que Tux en el ejemplo
    de Wikipedia \url{https://en.wikipedia.org/wiki/Block_cipher_mode_of_operation}.
    \item Realizar la operación con otro cifrador de 64 bits de bloque (des, des-ede3, etc) en modo ECB
    y comparar los resultados. El menor tamaño de bloque debería acentuar las diferencias en los
    cambios de color en la imagen original.
\end{enumerate}
